\documentclass[11pt]{amsart}

\usepackage{amsmath}
\usepackage{amssymb}
\usepackage{amsthm}
%\usepackage{psfig}

\begin{document}
\baselineskip=15pt
\textheight=8.4in
\parindent=0pt
\def\sk {\hskip .5cm}
\def\skv {\vskip .12cm}
\def\cos {\mbox{cos}}
\def\sin {\mbox{sin}}
\def\tan {\mbox{tan}}
\def\intl{\int\limits}
\def\lm{\lim\limits}
\newcommand{\frc}{\displaystyle\frac}
\def\xbf{{\mathbf x}}
\def\fbf{{\mathbf f}}
\def\gbf{{\mathbf g}}

\def\Ker{{\rm Ker\,}}
\def\phi{\varphi}

\def\dbA{{\mathbb A}}
\def\dbB{{\mathbb B}}
\def\dbC{{\mathbb C}}
\def\dbD{{\mathbb D}}
\def\dbE{{\mathbb E}}
\def\dbF{{\mathbb F}}
\def\dbG{{\mathbb G}}
\def\dbH{{\mathbb H}}
\def\dbI{{\mathbb I}}
\def\dbJ{{\mathbb J}}
\def\dbK{{\mathbb K}}
\def\dbL{{\mathbb L}}
\def\dbM{{\mathbb M}}
\def\dbN{{\mathbb N}}
\def\dbO{{\mathbb O}}
\def\dbP{{\mathbb P}}
\def\dbQ{{\mathbb Q}}
\def\dbR{{\mathbb R}}
\def\dbS{{\mathbb S}}
\def\dbT{{\mathbb T}}
\def\dbU{{\mathbb U}}
\def\dbV{{\mathbb V}}
\def\dbW{{\mathbb W}}
\def\dbX{{\mathbb X}}
\def\dbY{{\mathbb Y}}
\def\dbZ{{\mathbb Z}}


\def\char{{\rm char}}
\def\wt{{\rm wt}}
\def\Aut{{\rm Aut}}
\def\Gal{{\rm Gal}}
\def\Hom{{\rm Hom}}
\def\End{{\rm End}}
\def\Mat{{\rm Mat}}
\def\Irr{\mathrm {Irr}}
\def\Im{{\rm Im}}
\def\GL{{\rm GL}}
\def\SL{{\rm SL}}
\def\SNF{{\rm SNF}}
\def\tr {{\rm tr}}
\def\rk{{\rm rk}}
\def\lam{{\lambda}}

\def\la{{\langle}}
\def\ra{{\rangle}}

\bf
\begin{center}
{Homework \# 11, due on Canvas by 6pm on Sat, Apr 26th.}\rm
\end{center}
\vskip .2cm
\rm

{\bf Plan for the next 2 classes:} More on Galois correspondence, including further examples (14.4 in DF, Lecture 21 from Spring 2010 and Lectures 16-17 from Spring 2021). Cyclic extensions, including Hilbert's Theorem 90 and Kummer's Theorem (14.3 in DF, Lecture 23 from Spring 2010 and Lecture 19 from Spring 2021).
\vskip .1cm

%{\bf Note on hints:} Some hints are given at the end of the assignment, each on a separate page.
%Problems (or parts of problems) for which a hint is available at the end are marked with *.

\skv
{\bf Problem 1:} Let $p$ be a prime, $n$ a positive integer and
$\Phi_n(x)=x^{p^n}-x\in \mathbb F_p[x]$. Prove that $\Phi_n(x)$ is equal to the product
of all monic irreducible polynomials in $\mathbb F_p[x]$ whose degrees divide $n$
(where each polynomial occurs with multiplicity one). {\bf Hint:} Reformulate the question in terms of $\overline{\mathbb F_p}$-roots
of $\Phi_n(x)$.
\skv

{\bf Problem 2:} Let $K\subset\dbC$ be the splitting field of $f(x)=x^4-2$ over $\dbQ$.

\begin{itemize}
\item[(a)] Choose an order on the set of roots of $x^4-2$ and describe the associated embedding
of $\Gal(K/\dbQ)$ to $S_4$.

\item[(b)] Describe all subgroups of $\Gal(K/\dbQ)$ and the corresponding subfields of $K$ (with proof).

\end{itemize}

\skv
{\bf Problem 3:} Let $p$ and $q$ be distinct primes with $q>p$,
and let $K/F$ be a Galois extension of degree $pq$. Prove that
\begin{itemize}
\item[(a)] There exists a field $L$ with $F\subseteq L\subseteq K$ and $[L:F]=q$

\item[(b)] There exists a unique field $M$ with $F\subseteq M\subseteq K$ and $[M:F]=p$.
\end{itemize}
\skv

\skv
{\bf Problem 4:} Let $K/F$ and $L/F$ be Galois extensions.

\begin{itemize}
\item[(a)] Prove that the extension $KL/F$ is also Galois and there is a natural
embedding $\iota:\Gal(KL/F)\to \Gal(K/F)\times \Gal(L/F)$.

\item[(b)]
Assume now that $K/F$ and $L/F$ are both finite. Prove that the map $\iota$ in (a)
is an isomorphism if and only if $K\cap L=F$.
\end{itemize}

{\bf Problem 5:} \rm Before doing this problem, read the first half
of Section 14.4 in DF (pp. 591-593).

{\bf Definition~1:} Let $L/F$ be a finite separable extension
and let $\overline F$ be an algebraic closure of $F$ containing $L$.
A subfield $L'$ of $\overline F$ is called
{\bf conjugate to $L$ over $F$} if $L'=\sigma(L)$ for some $F$-embedding of
$\sigma$ into $\overline F$. Note that $L/F$ is Galois if and
only if $L$ does not have any $F$-conjugates besides $L$ itself.
\skv
{\bf Definition~2:} A finite extension $K/F$ is called
a {\bf $p$-extension} if $K/F$ is Galois and $Gal(K/F)$ is a $p$-group.
\skv
\begin{itemize}
\item[(a)] Let $L/F$ be a separable extension of degree $n$, and
let $K$ be the Galois closure of $L$ over $F$ (see p. 594 in DF or subsection 19.2 in Lecture~19 from Spring~2010 for the definition). 
Prove that $K$ can be written as a compositum $L_1L_2\ldots L_n$
where $L_1,\ldots L_n$ are (not necessarily distinct) conjugates of
$L$ over $F$.

\item[(b)] Let $K/F$ and $L/F$ be finite $p$-extensions. Prove
that $KL/F$ is also a $p$-extension.

\item[(c)] Suppose $K/L$ and $L/F$ are both $p$-extensions, and
let $M$ be the Galois closure of $K$ over $F$ (note: we do
not know whether $K/F$ is Galois or not). Prove that
$M/F$ is also a $p$-extension. {\bf Hint:} first use (b) to show that
$M/L$ is a $p$-extension.

\item[(d)] Now assume only that $L/F$ is a separable extension
with $[L:F]$ a power of $p$, and let $M$ be the Galois closure
of $L$ over $F$. Prove that $[M:F]$ need not be a power of $p$.
\end{itemize}
\skv
{\bf Problem 6:} Let $p$ be a prime, $K=\mathbb F_p(t)$, the field of rational functions in one variable over $\dbF_p$.
Let $\sigma:K\to K$ be the unique automorphism of $K$ such that $\sigma(t)=t+1$ and $G=\langle \sigma\rangle$ (clearly
$G$ is cyclic of order $p$). Let $F=K^G$. Find an explicit element $s$ such that $F=\dbF_p(s)$ and prove your answer.
\skv
{\bf Problem 7:} 
Let $E/F$ be a field extension, let $K_1$ and $K_2$ be subfields of $E$ containing $F$,
and assume that the extensions $K_1/F$ and $K_2/F$ are finite. Let $K_1 K_2$ be the composite
field of $K_1$ and $K_2$. Prove that the $F$-algebra $K_1\otimes_F K_2$ is a field if and only if
$$[K_1 K_2: F]=[K_1: F][K_2: F].$$ {\bf Hint:} Use the universal property of tensor products to relate
$K_1\otimes_F K_2$ and $K_1 K_2$.
\skv
{\bf Problem 8:} Let $K/F$ be a finite Galois extension of degree $n$. Prove that
$K\otimes_F K\cong \underbrace{K\times\ldots\times K}_{n \mbox{ times }}$ as $K$-algebras. {\bf Hint:} Use the
Primitive Element Theorem to represent the second $K$ in $K\otimes_F K$ in a suitable way, so that the tensor
product can be computed using techniques we discussed at the start of the semester.
\end{document}
